\documentclass[11pt,a4paper]{article}
\usepackage{fontspec}
\setmainfont{Noto Serif}
\setsansfont{Noto Sans}
\usepackage{babel}
\babelprovide[main,import]{russian}
\usepackage{amsmath,amssymb,mathtools,booktabs,geometry,fancyhdr,microtype,xcolor,hyperref}
\geometry{left=21mm,right=21mm,top=21mm,bottom=21mm}
\definecolor{ink}{HTML}{173D57}
\hypersetup{colorlinks=true,urlcolor=ink,linkcolor=ink,pdftitle={BGS-F. Четвёртый момент и знаковая резольвента}}
\pagestyle{fancy}\fancyhf{}\lhead{\small BGS-F}\rhead{\small Четвёртый момент и знаковая резольвента}\cfoot{\thepage}
\setlength{\headheight}{15pt}\setlength{\parindent}{1em}\setlength{\parskip}{2.5pt}
\setlength{\emergencystretch}{3em}\tolerance=3000
\newcommand{\Tr}{\operatorname{Tr}}
\newcommand{\EE}{\mathbb E}
\newcommand{\R}{\mathbb R}
\newcommand{\Z}{\mathbb Z}
\newcommand{\C}{\mathbb C}
\newcommand{\si}{\mathrm{Sine}_2}
\newcommand{\eps}{\varepsilon}
\newcommand{\sinc}{\operatorname{sinc}}
\newcommand{\sect}[1]{\section*{\sffamily\color{ink}#1}}
\newcommand{\subt}[1]{\par\medskip\noindent\textbf{#1.}\quad}
\newcommand{\T}{\mathbb T}
\newcommand{\dist}{\operatorname{dist}}
\newcommand{\Ad}{\operatorname{Ad}}
\newcommand{\op}{\mathrm{op}}
\newcommand{\HS}{\mathrm{HS}}
\newcommand{\pa}{\mathrm p}
\newcommand{\Cyc}{\mathcal C}
\begin{document}
\noindent{\small\sffamily\color{ink} МАТЕМАТИЧЕСКАЯ ФИЗИКА\hfill 20 НОЯБРЯ 2026}
\begin{center}{\LARGE\sffamily Четвёртый момент и знаковая резольвента}\\[4pt]
{\large Спектральные кластеры почти до границы масштаба \(N^{-3/2}\)}
\end{center}
Пусть \(N=7^k\), \(m=3N/7\) и \(U_N(t)=e^{itA_N}M_N\) для исходных \(A_N,M_N\). Физическая сила равна \(u_N=t/N\). Размеры спектральных щелей измеряются после развёртки, у которой средняя щель равна единице.

\subt{Локальная теорема}
При
\[
             u_N=o\!\left(N^{-3/2}(\log N)^{-1/4}\right)
\]
весь локальный процесс сходится к случайно сдвинутой решётке шага \(7/3\) с кратностями \(2(1+\nu_7(B+j))\), где \(B\) равномерен в \(\mathbb Z_7\). Для каждого фиксированного \(T<\infty\)
\[
     \sup_{|r|\le TN}\frac{|\Tr U_N(t)^r-\Tr M_N^r|}{N}\longrightarrow0.
\]
Это, в частности, даёт когерентный обычный след при \(r=3N/7\). При малой ненулевой константе перед \(N^{-3/2}(\log N)^{-1/4}\) предельный синусный закон также исключается. Предыдущая оценка полной резольвенты требовала множитель \((\log N)^{-1}\).

\subt{Парная теорема}
Для каждой фиксированной тестовой функции Шварца парные статистики линейного и возмущённого спектров имеют одну и ту же асимптотику, если
\[
              u_N=o\!\left(N^{-8/5}(\log N)^{-2/5}\right).
\]
В частности, предел полного исходного парного теста \(\Delta_{2,N}(b_*)\) принадлежит \((0.08443,0.08656)\), как у линейного оператора.

\subt{Основная оценка}
Для каждого фиксированного \(R\in(0,1]\) построены непересекающиеся окружности радиуса \(R/m\) вокруг линейных фаз. При \(|t|\le N^{-1/2}k^{-1/4}\) исходная кратность сохраняется внутри всех, кроме доли
\[
                          O_R(t^4 N^2 k)
\]
окружностей. Доказательство использует точное уравнение для замкнутого слова из четырёх вейлевых меток, четвёртые нормы Шаттена и знаки дальнего резольвентного блока. При каждом конечном \(k\) все метки \(A_N\) берутся из исходной модели. Предел \(\mathrm{Sine}_2\) при фиксированном \(u_N=10^{-80}\) остаётся открытым, поскольку эта сила выходит из указанной области при \(k\to\infty\).
\newpage
\sect{1. Спектральные веса и их преобразование Фурье}
Нормируем \(M_N\) делением на общую фазу \(\lambda_0\), так что \(V=M_N/\lambda_0\), \(V^m=I\). Пусть \(E_j\) --- спектральный проектор на \(e^{2\pi i j/m}\) для \(V\). Кратность \(d_j\) равна
\[
           2\sum_{a=0}^{k-1}\mathbf1_{7^a\mid j}+\mathbf1_{j=0}.
\]
В частности, \(d_j\le 2k+1\) и \(\sum_j d_j^2=(8N-5)/3\). Эти формулы следуют из двух циклов длины \(3\cdot7^{a-1}\) на каждом \(7\)-адическом слое аддитивного порядка \(7^a\). Они подробно выведены в части о резольвенте.

Как в той части, для \(D=(m-1)/2\) задаём чётные веса на \(\Z/m\Z\) по представителям минимального модуля:
\[
 w_0=R^{-1},\qquad
 w_d=(2m\sin(\pi |d|/m)-R)^{-1}\quad(1\le |d|\le D).
\]
На окружности \(\Gamma_j=\{z:|z-\lambda_0 e^{2\pi i j/m}|=R/m\}\)
\(|\lambda_i-z|^{-1}\le m w_{i-j}\). Для неотрицательного веса \(g\) пишем \(Z_j(g)=\sum_i\sqrt{g_{i-j}}E_i\). При \(1\le L<D\) определим \(w^{\rm n}=w\mathbf1_{|d|\le L}\), \(w^{\rm f}=w\mathbf1_{|d|>L}\). Преобразование Фурье \(\widehat g(h)=\sum_d g_d e^{-2\pi i h d/m}\) берётся без нормировки.

Для \(1\le d\le D\) оценка из части о резольвенте даёт
\(0\le w_d-1/(2\pi d)\le 3d/(5m^2)+R/(18d^2)\). Правая часть суммируется с общей границей \(O_R(1)\). Отсюда для \(g=w^{\rm n}\) либо \(g=w\)
\begin{equation}
 \begin{split}
 \sup_h|\widehat w^{\rm n}(h)|&\le C_R(1+\log(2L)),\\
 |\widehat g(h)|&\le C_R\bigl(1+\log(m/(1+|h|_m))\bigr),\quad h\ne0 ,
 \end{split}                                                    \tag{1}
\end{equation}
где \(|h|_m\) --- циклическое расстояние до нуля. Для второго неравенства сумму \(\sum_{d\le L}e^{2\pi i hd/m}/d\) разбиваем при \(d\) порядка \(m/|h|_m\). Перед разбиением используем гармоническую сумму, после него --- суммирование по частям и ограничение геометрических сумм через \(m/|h|_m\). Для полного \(w\) первое неравенство действует с \(L=D\).

Дальний вес монотонно убывает на \([L+1,D]\), поэтому суммирование по частям даёт
\begin{equation}
 |\widehat w^{\rm f}(h)|\le C_R\min\{k,\ m/(L|h|_m)\},\quad h\ne0,
 \qquad |\widehat w^{\rm f}(0)|\le C_R k .                       \tag{2}
\end{equation}
На подгруппе кратных \(q\mid m\) отсюда следуют две границы для суммы квадратов по ненулевым \(h\):
\begin{equation}
 \sum_{\substack{q\mid h\\h\ne0}}|\widehat g(h)|^2\le C_Rm/q
             \quad(g=w^{\rm n},w),\qquad
 \sum_{\substack{q\mid h\\h\ne0}}|\widehat w^{\rm f}(h)|^2
                           \le C_R (m/q)k/L.                   \tag{3}
\end{equation}
Для первой сравниваем сумму с интегралом \((1+\log(1/x))^2\) на \((0,1)\). Для второй используем \(\sum_{j\ge1}\min(k^2,(b/j)^2)\le2bk\) при \(b=m/(qL)\), по интегральному сравнению.
\newpage
\sect{2. Замкнутое слово из четырёх одинаковых меток}
В \(\mathbb Z_7\) выберем \(\rho^2=2,\ \rho\equiv3\bmod7\). Симплектические координаты \((a+\rho b,(a-\rho b)/(-2\rho))\) диагонализуют действие \(S\). Его диагональные множители равны \(\gamma=3+2\rho\) и \(\gamma^{-1}\). У каждой метки носителя \(A_N\) обе новые координаты имеют один порядок \(d\in\{0,1\}\), а после деления на \(7^d\) являются единицами. Совпадение вейлевых меток тогда проверяется по модулю \(7^K\), \(K=k-d\). Порядок \(\gamma\bmod7^K\) равен \(o=3\cdot7^{K-1}\) и \(m/o\le7\).

Рассмотрим ненормированную норму Шаттена \(S_4\) матрицы \(mZ_j(g)T(v)Z_j(g)\) с одной вейлевой меткой \(v\). Её четвёртая степень выражается следом четырёх операторов \(T(v),T(-v),T(v),T(-v)\), разделённых четырьмя одинаковыми весами \(\sum_i g_{i-j}E_i\). Разложение каждого веса по степеням \(V\) вносит коэффициент \(\widehat g(h)e^{-2\pi i hj/m}/m\). Усреднение по \(j\) оставляет только четыре частоты с суммой ноль, а множитель \(m^4\) сокращает четыре знаменателя.

После переноса степеней \(V\) вправо по точной ковариантности Вейля времена меток равны \(0,h_1,h_1+h_2,h_1+h_2+h_3\). Если соответствующие три степени \(\gamma\) обозначены \(x,y,z\), условие ненулевого следа имеет вид
\[
             1-x+y-z=0,\qquad 1-x^{-1}+y^{-1}-z^{-1}=0
                         \quad\bmod 7^K.
\]
Поскольку \(1+y\) --- единица (по модулю 7 группа \(\langle\gamma\rangle\) состоит из \(1,2,4\)), имеем \(y=xz\) и
\begin{equation}
                        (x-1)(z-1)=0\bmod7^K .                  \tag{4}
\end{equation}
Обратно, для любых \(x,z\) из \(\langle\gamma\rangle\) с (4) значение \(y=xz\) даёт замкнутое слово. Дополнительная фаза Вейля имеет модуль один, поэтому след такого слова по модулю \(N\).

Положим \(\nu(a)=\min(K,\nu_7(\gamma^a-1))\), \(a\bmod o\), и
\(G(a)=\sum_{h\equiv a\bmod o}|\widehat g(h)|\). Из \(y=xz\) следует, что четыре частоты по модулю \(o\) имеют вид \(a,b,-a,-b\), где \(x=\gamma^a,z=\gamma^b\). Отбрасывая условие связи при поднятии всех четырёх частот по модулю \(m\) и используя чётность \(G\), получаем
\begin{equation}
 \EE_j\|mZ_j(g)T(v)Z_j(g)\|_{S_4}^4
     \le N\sum_{\substack{a,b\bmod o\\\nu(a)+\nu(b)\ge K}}G(a)^2G(b)^2 .
                                                                 \tag{5}
\end{equation}

\subt{Точный подсчёт без весов}
Для \(s\ge1\) условие \(\nu(a)\ge s\) задаёт кратные \(3\cdot7^{s-1}\). При ненулевом \(a\) порядка \(s<K\) число подходящих \(b\) равно \(7^s\), включая \(s=0\). При \(a=0\) оно равно \(o\). Всего \((6K-1)7^{K-1}\) слов с первым временем \(0\). Поэтому несокращённый счёт четырёх меток действительно имеет линейный множитель \(K\), даже с одной парой противоположных исходных частот.
\newpage
\sect{3. Четвёртые нормы ближнего и дальнего блоков}
По Коши \(G(a)^2\le(m/o)\sum_{h\equiv a\bmod o}|\widehat g(h)|^2\). Для \(a\ne0,\ \nu(a)\ge s\) нужные \(h\) лежат в подгруппе кратных \(3\cdot7^{s-1}\), а при \(s=0\) ограничение отсутствует. Поэтому (3) даёт
\[
 \sum_{\substack{a\ne0\\\nu(a)\ge s}} G(a)^2
 \le C_R\,m\,7^{-s}
           \begin{cases}1,& g=w^{\rm n},w,\\ k/L,&g=w^{\rm f}.\end{cases}
                                                               \tag{6}
\]
При \(a=0\) имеем соответственно \(G(0)\le C_R(1+\log(2L))\) для ближнего веса и \(G(0)\le C_R k\) для полного и дальнего.

Сначала используем \(g=w^{\rm n}\). В (5) пары с \(a=0\) или \(b=0\) дают \(C_Rm(1+\log(2L))^2\) до множителя \(N\). Для остальных при \(\nu(a)=s\) умножаем (6) для \(s\) и \(K-s\), получая \(C_R m^2 7^{-K}\le C_R m\) для каждого \(s\). Значит,
\[
    \EE_j\|mZ_j(w^{\rm n})T(v)Z_j(w^{\rm n})\|_{S_4}^4
                      \le C_R N^2 \bigl[k+(1+\log(2L))^2\bigr].
\]
Для полного \(w\) тот же довод даёт \(C_R N^2 k^2\). Для дальнего \(w^{\rm f}\) из (6) получаем \(C_R m k^3/L\) от пар с нулём и \(C_R m k^3/L^2\) от остальных. В первую границу включается \(G(0)^4=O_R(k^4)\), если \(L\le m/k\). Именно этот диапазон применяется далее. В результате дальняя четвёртая степень не больше \(C_R N^2 k^3/L\).

Неравенство треугольника в норме \(S_4\) и затем в \(L^4(j)\) позволяет подставить \(A_N=\sum_v c_v T(v)\) вместо \(T(v)\). Число меток конечно и \(\sum_v|c_v|\le593/500\), поэтому меняется лишь постоянная. Выберем отныне \(L=\lfloor k^{16}\rfloor\), что для больших \(k\) удовлетворяет \(7<L<D\) и \(L\le m/k\). Пишем \(Z_{\rm n}, Z_{\rm f}, Z\) для соответствующих весов с центром \(j\). При такой отсечке \((1+\log(2L))^2=O(k)\). Получаем три оценки
\begin{equation}
 \begin{aligned}
  \EE_j\|tmZ_{\rm n} A_N Z_{\rm n}\|_{S_4}^4 &\le C_R t^4 N^2 k,\\
  \EE_j\|tmZ_{\rm f} A_N Z_{\rm f}\|_{S_4}^4 &\le C_R t^4 N^2 k^3/L,\\
  \EE_j\|tmZ A_N Z\|_{S_4}^4 &\le C_R t^4 N^2 k^2 .
 \end{aligned}                                                   \tag{7}
\end{equation}
В частности, третья строка ограничивает и четвёртую степень операторной нормы смешанного блока \(tmZ_{\rm n}A_N Z_{\rm f}\), по сжатию на две области.

Использование только третьей строки уже укорачивает логарифмическую потерю до \((\log N)^{1/2}\) в условии на \(t\). Для дальнейшего улучшения требуется учитывать фазу резольвенты в квадратичном блоке.
\newpage
\sect{4. Дальняя резольвента с фазой}
Пусть \(P_{\rm f}(j)=\sum_{|i-j|_m>L}E_i\) и
\[
     F_j=\frac1m P_{\rm f}(j)M_N(M_N-\lambda_jI)^{-1}\big|_{\rm f}
       =\sum_i f_{i-j}E_i,\qquad
     f_d=\frac{\mathbf1_{|d|_m>L}}{m(1-e^{-2\pi i d/m})}.
\]
Выражение с обратным понимается только на указанном дальнем подпространстве и продолжено нулём. В \(\widehat f\) сохраняются исходные знаки: для \(1\le d\le D\)
\[
       \frac1{m(1-e^{-2\pi i d/m})}
         =\frac1{2m}-\frac{i}{2m}\cot(\pi d/m)
         =\frac1{2m}-\frac{i}{2\pi d}+O(d/m^2).
\]
Для отрицательного \(d\) меняется знак мнимой части. Ошибка равномерна, поскольку \((1/x-\cot x)/x\) ограничено на \(0<x\le\pi/2\). Вещественная часть и ошибка суммируются в \(O(1)\). Главный мнимый член после сложения положительных и отрицательных \(d\) есть усечённая синусная сумма. Она ограничена равномерно по частоте и по обеим границам отрезка: при \(\theta=2\pi |h|_m/m\) члены с \(d\le1/\theta\) ограничиваем через \(|\sin(\theta d)|/d\le\theta\), после чего применяем суммирование по частям к геометрическим суммам. При \(h=0\) синусный вклад нулевой. Таким образом,
\begin{equation}
                              \sup_h|\widehat f(h)|\le C .      \tag{8}
\end{equation}

Нужно оценить \(Z_{\rm n} A_N F_j A_N Z_{\rm n}\) в норме Гильберта--Шмидта. По треугольнику в пространстве \(L^2(j;\mathrm{HS})\) достаточно рассмотреть
\(Z_{\rm n} T(v) F_j T(w) Z_{\rm n}\) для каждой пары исходных меток \(v,w\). Квадрат нормы выражается через четыре метки \(v,w,-w,-v\), разделённые весами \(F_j,\ Z_{\rm n}^2,\ F_j^*,\ Z_{\rm n}^2\) в этом циклическом порядке.

Раскладываем веса по степеням \(V\) и усредняем по \(j\), как в разделе 2. Множитель \(m^4\) сокращает четыре нормировки. Пусть времена меток суть \(0,s_1,s_2,s_3\) и \(X=\gamma^{s_3},Y=\gamma^{s_1},Z_0=\gamma^{s_2}\). Тогда по (8) произведение модулей четырёх частотных весов не больше
\[
       C|\widehat w^{\rm n}(s_2-s_1)|\,|\widehat w^{\rm n}(-s_3)|.
\]
Оно ограничивается полусуммой квадратов двух последних множителей с постоянной. Для квадрата, зависящего от \(s_2-s_1\), циклически начинаем слово с метки \(-w\) и снимаем общее время \(s_2\). Роли \(v,w\) меняются с общим изменением знака, а задача сводится к квадрату веса от заключительного времени. Остаётся подсчитать для каждого \(s_3\) число завершений \(s_1,s_2\) замкнутого слова.
\newpage
\sect{5. Подсчёт смешанных меток}
Снимем общую степень семи с обеих меток. По модулю \(7^K\), \(K=k\) или \(k-1\), их координаты можно записать \(v=7^{\delta_0}(a,b)\), \(w=7^{\delta_1}(c,d)\), где \(a,b,c,d\) обратимы, \(\delta_i\in\{0,1\}\) и \(\min(\delta_0,\delta_1)=0\). Первое уравнение замыкания и эквивалентный вид второго после использования первого имеют вид
\begin{equation}
  7^{\delta_1} c(Y-Z_0)= 7^{\delta_0}a(X-1), \qquad
  7^{\delta_0}(X-1)(YZ_0-\eta X)=0,\quad \eta=\frac{ad}{bc}.
                                                               \tag{9}
\end{equation}
Для второго равенства домножаем уравнение
\(7^{\delta_0}b(1-X^{-1})+7^{\delta_1}d(Y^{-1}-Z_0^{-1})=0\)
на \(XYZ_0\) и подставляем первое. Число поднятий каждого из времён из периода \(3\cdot7^{K-1}\) к \(m\) не больше семи.

Пусть \(s=\min(K,\nu_7(X-1))\). При \(s+\delta_0\ge K\) таких \(X\) не более семи, и с каждым из них число завершений \(O(m)\) по первому равенству. Этот вклад после веса \(|\widehat w^{\rm n}|^2\) не больше \(C_R m(1+\log(2L))^2\). При \(s+\delta_0<\delta_1\) решений нет. В остальных случаях положим
\[
       K'=K-s-\delta_0,\qquad C_0=(a/c)7^{\delta_0-\delta_1}(X-1).
\]
Первое равенство определяет \(Z_0=Y-C_0\bmod7^{K-\delta_1}\); второе эквивалентно \(Y^2-C_0Y-\eta X=0\bmod7^{K'}\). Здесь \(1\le K'\le K-\delta_1\).

Если \(s+\delta_0>\delta_1\), дискриминант \(C_0^2+4\eta X\) обратим по модулю семь, так что у квадратного сравнения не более двух корней. Число завершений на данном \(s\) не больше \(C7^s\) после учёта поднятий. Ненулевые веса по временам \(s_3\) с данным порядком лежат в подгруппе частот, где по (3) их сумма квадратов \(\le C_Rm7^{-s}\). Получаем \(C_Rm\) на каждый порядок \(s\), всего \(C_Rmk\).

Остались случаи \(s+\delta_0=\delta_1\). Если \(s=\delta_0=\delta_1=0\), дискриминант по \(X\) имеет вид \((a/c)^2(X-1)^2+4\eta X\). Его корни по модулю семь простые, кроме возможности \(\eta=(a/c)^2\bmod7\), когда единственный потенциальный корень \(-1\) не лежит в группе. Если \((\delta_0,\delta_1,s)=(0,1,1)\), пишем \(X=1+7X_0\). Дискриминант по \(X_0\) имеет простые корни по модулю семь, поскольку \(\eta\) обратима. Гензелево поднятие в обоих случаях даёт хвост \(O(m7^{-b_0})\) для числа \(s_3\), где дискриминант делится на \(7^{b_0}\), \(b_0\le K'\); во втором случае считаем \(X\) по модулю \(7^{b_0+1}\). При порядке дискриминанта \(b_0\), усечённом на \(K'\), число корней \(\le 2\cdot7^{b_0/2}\) по сведению к квадратному корню и подъём \(7^{s+\delta_0+\delta_1}\) ограничен постоянной. Сумма по \(b_0\) есть \(O(m)\) без веса. Домножаем на верхнюю границу \(|\widehat w^{\rm n}|^2\) из (1).

Итак, взвешенное число завершений не больше \(C_R m[k+(1+\log(2L))^2]\). След каждого замкнутого слова ограничен \(N\). Возвращая множители и все коэффициенты \(A_N\), получаем для \(L=\lfloor k^{16}\rfloor\)
\begin{equation}
 \EE_j\|t^2m^2 Z_{\rm n} A_N F_j A_N Z_{\rm n}\|_{\HS}^2
                           \le C_R t^4 N^2 k.                  \tag{10}
\end{equation}
\newpage
\sect{6. Сохранение кратностей через дополнение Шура}
На \(\Gamma_j\) положим \(D_z=(M_N-zI)^{-1/2}\) и
\(B(z)=D_z\Delta D_z\), \(\Delta=(e^{itA_N}-I)M_N\). Для проверки числа корней достаточно обратимости \(I+s B(z)\) для \(0\le s\le1\) на всём контуре. Обозначим через \(B^{(1)}\) линейный блок \(itD_z A_N M_N D_z\). Все нормы этого блока на \(\Gamma_j\) покомпонентно ограничены нормами соответствующих блоков \(tm Z A_N Z\), так как \(|D_z|\le \sqrt m Z\) в диагональном базисе. Смешанные нормы линейного блока ограничены одним числом \(Y_j=\|tmZ_{\rm n}A_N Z_{\rm f}\|_{\op}\).

Пусть \(|t|\le N^{-1/2}k^{-1/4}\). Остаток \(B-B^{(1)}\) равномерно по контуру имеет нормы \(O_R(t^2 N)\) в ближнем, \(O_R(t^2N/L)\) в дальнем и \(O_R(t^2N/\sqrt L)\) в смешанных блоках. Это следует из \(\|e^{itA_N}-I-itA_N\|_{\op}\le t^2\|A_N\|_{\op}^2/2\) и максимальных весов \(R^{-1}\) и \(O(1/L)\). В частности, ближний остаток \(o(1)\).

Выбираем хорошие узлы условиями
\[
 \begin{gathered}
 \|tmZ_{\rm n} A_N Z_{\rm n}\|_{\op}\le1/4,\quad
 \|t^2m^2 Z_{\rm n} A_N F_j A_N Z_{\rm n}\|_{\op}\le1/4,\\
 \|tmZ_{\rm f} A_N Z_{\rm f}\|_{\op}\le k^{-2},
                         \qquad Y_j\le k^{3/4}.
 \end{gathered}
\]
По (7),(10) и Маркову доля остальных не больше
\begin{equation}
     \beta\le C_R t^4 N^2\bigl(k+k^{11}/L+k^{-1}\bigr)
                        \le C_R' t^4 N^2 k .                   \tag{11}
\end{equation}
Для последнего условия использована норма полного блока из третьей строки (7), возведённая в четвёртую степень. Для дальнего условия порог даёт множитель \(k^8\).

У хорошего узла дальний блок \(B_{\rm ff}\) имеет норму \(O(k^{-2})\). Устраняем \(I+sB_{\rm ff}\) дополнением Шура. Замена \((I+sB_{\rm ff})^{-1}\) на \(I\) и смешанных блоков на их линейные части даёт в ближнем дополнении общий остаток
\(O_R(k^{-2}k^{3/2}+k^{3/4}/\sqrt L+1/L)=o(1)\).
Остаётся квадратичный блок \(B_{\rm nf}^{(1)}(z)B_{\rm fn}^{(1)}(z)\). Его внутренний множитель есть дальняя часть \(M_N(M_N-zI)^{-1}\). От центра \(\lambda_j\) до контура разность этого множителя покомпонентно \(\le m w_{i-j}R/(4L)\): используем тождество резольвент и \(|\lambda_i-\lambda_j|\ge4|i-j|_m/m\). Замена на центральный множитель \(m F_j\) поэтому стоит не больше \(C_R Y_j^2/L=o(1)\). Оставшийся блок ограничен \(1/4\) по второму условию. Таким образом, ближнее дополнение отклоняется от \(I\) не больше чем на \(1/2+o(1)\), равномерно по всему \(\Gamma_j\) и \(s\in[0,1]\). Оба блока обратимы.

Число собственных значений внутри \(\Gamma_j\) постоянно на пути \(M_N+s\Delta\) по принципу аргумента. Следовательно, при \(U_N(t)\) оно точно равно \(d_j\). Это завершает доказательство оценки доли (11), включая нелинейный остаток исходной экспоненты.
\newpage
\sect{7. Предельные выводы}
Сопоставим все \(d_j\) фаз в хорошем круге с его центром. Из формулы кратностей доля \(\epsilon\) несопоставленных фаз удовлетворяет
\[
 \epsilon\le \frac67\sum_{a=0}^{k-1}\min(\beta,7^{-a})+\frac1N,
       \qquad \epsilon\le C k\beta.
\]
Первое выражение не больше \(C\beta(1+\log(1/\beta))+1/N\). Каждая сопоставленная фаза удалена от центра на \(O(R/N)\), а прочие уровни имеют суммарную массу \(\epsilon\).

Если \(t=o(N^{-1/2}k^{-1/4})\), то (11) даёт \(\epsilon\to0\) при любом фиксированном \(R>0\). Усреднение по общему равномерному началу развёртки даёт для любого \(f\in C_c^1(\R)\)
\[
 \EE|\Xi_N^{U}(f)-\Xi_N^{M}(f)|
                     \le C R\|f'\|_1+2\epsilon\|f\|_1 .
\]
Переходя сначала к \(k\to\infty\), затем к \(R\downarrow0\), переносим локальный закон линейного спектра. Он получается непосредственно из \(d_j\): непрерывный сдвиг равномерен по периоду \(7/3\), а соседние кратности определяет общий равномерный \(B\in\mathbb Z_7\). Тем же сопоставлением \(\sup_{|r|\le TN}N^{-1}|\Tr U_N^r-\Tr M_N^r|\le C_T R+2\epsilon\), что доказывает утверждение об обычных следах. При \(|t|\le c N^{-1/2}k^{-1/4}\) оценка доли равна \(C_Rc^4\). Фиксируем достаточно малый \(R\), затем выбираем \(c>0\) достаточно малым. Ограниченный локальный тест на две близкие точки из части о резольвенте тогда даёт положительный разрыв с \(\si\).

Для \(\psi\in\mathcal S(\R)\) возьмём парную сумму \(P_\psi(V)=N^{-1}\sum_{a\ne b}\sum_{\ell\in\Z}\psi(N(\theta_a-\theta_b+\ell))\). Доказанная ранее оценка для такого сопоставления имеет вид
\[
       |P_\psi(U_N(t))-P_\psi(M_N)|
                \le C_\psi\bigl(R+\epsilon(1+N|t|+k)\bigr).
\]
Для сопоставленных пар она следует из теоремы о среднем и \(\sum_j d_j^2=O(N)\). Для остальных пар используется равномерная граница \(C(1+N|t|+k)\) на число фаз в интервале единичной развёрнутой длины. Последняя следует из скорости фаз не выше \(\|A_N\|_{\op}\) в радианах по параметру \(t\) и подсчёта исходных кратностей на дуге. Теперь по (11)
\[
            \epsilon(1+N|t|+k)\le
                      C_R t^4N^2 k^2(1+N|t|+k).
\]
Это стремится к нулю при \(t=o(N^{-3/5}k^{-2/5})\). Снова \(R\downarrow0\) переносит все парные тесты. Подстановка \(\widehat\psi=|\widehat b_*|^2\) даёт перенос полного \(\Delta_{2,N}(b_*)\) и его положительного предела. Формулы для физической силы получаются делением \(t\) на \(N\); \(k=(\log N)/\log7\).

\subt{Проверка}
Целочисленный код проверяет (4) и точное число замкнутых слов, а также смешанное сведение (9) для всех пар исходных меток при \(N=49\) и выборочных пар при \(N=343\). Независимый матричный скрипт сравнивает обе формулы четвёртого порядка с прямыми вычислениями на \(N=49\). Асимптотика при фиксированном \(10^{-80}\) из (11) не следует.
\end{document}
